% ===================================================================
% Track 1 of Sprint 1 dictionary-completion: Phillips-Kleinman /
% core-valence orthogonality projection.
%
% Drafted: 2026-05-15
% Target insertion point: Paper 34, §III, immediately after §III.19
% (nuclear tensor multipole, label sec:proj_tensor_multipole) and
% before the §III "Boundary of \S\ref{sec:layer2}" subsection at
% line 816 of papers/observations/paper_34_projection_taxonomy.tex.
%
% ===================================================================
% §IV three-axis table row (to be inserted in tab:projection_axes
% between the tensor-multipole row and the \bottomrule):
%
%   Phillips--Kleinman orthogonality & core orbital indices
%     $\{\phi_c\}$ + binding energies $\{E_c\}$ & dimensionless
%     (orthogonality is dimensionless) &
%     ring-preserving over $\mathbb{Q}(\alpha)$; no transcendental
%     injected beyond what the source spectrum already carries
%
% One-sentence summary for the table caption / changelog:
%   "Phillips--Kleinman / core-valence orthogonality is a Layer-2
%   $\to$ Layer-2 operator-level projection that takes a valence-only
%   Hamiltonian and adds the pseudopotential / projection content
%   that enforces core-valence orthogonality, with the projected-out
%   core orbitals as external Layer-2 input (Clementi-Roetti HF or
%   the framework's own core-shell solution)."
%
% Cross-references added in this draft:
%   - Paper 17 §IV (sec:pk, existing same-center PK treatment)
%   - Paper 17 §VI (sec:limitations / outlook, l_max divergence;
%     channel-blind vs $l$-dependent PK)
%   - Paper 19 §sec:w1c_residual (W1c-residual wall; PK cross-center
%     engineering closure with honest negative on NaH binding)
%   - CLAUDE.md §3 row "Phillips-Kleinman cross-center barrier for
%     second-row chemistry binding (post-Track-2 sprint, 2026-05-08)"
%   - Track CB negative result (coupled composition: cross-block ERIs
%     replacing PK), CLAUDE.md §3 row "Coupled composition (cross-block
%     ERIs replacing PK)"
%   - §III.17 nuclear charge-density (sibling Layer-2 input class)
%   - §III.18 nuclear magnetization-density (sibling Layer-2 input class)
%   - §III.19 nuclear tensor multipole (sibling Layer-2 input class)
%
% Suggested §V.B (off-precision matches catalogue) row, *if* the PM
% chooses to add one in the integration pass:
%
%   NaH balanced FCI (n_max=2) | PK cross-center + W1c screening |
%   E_min at smallest tested R, no internal equilibrium minimum |
%   error code S (structural floor: valence basis on frozen-core
%   center is Z_orb=1 hydrogenic, not screened Schrödinger
%   eigenstates of Z_eff(r)) | Paper 19 sec:w1c_residual;
%   CLAUDE.md §2 PK cross-center sprint bullet; CLAUDE.md §3 row.
%
% No new bibliography entries needed.  The existing bibitem \texttt{pk1959}
% (Phillips--Kleinman 1959) is already in Paper 34's bibliography (via
% the Paper 17 cross-reference chain).  Clementi--Roetti is cited in
% Paper 17 / Paper 19 already.
% ===================================================================

\subsection{Phillips--Kleinman / core--valence orthogonality}
\label{sec:proj_phillips_kleinman}
\textbf{Source:} a valence-only Hamiltonian on the Fock-projected $S^3$
graph in which the core orbitals $\{\phi_c\}$ of a heavy atom are
\emph{frozen out} of the active Hilbert space, leaving the valence
electrons without any structural memory of the Pauli-exclusion
constraint they must satisfy with respect to the core.\\
\textbf{Target:} same valence-only Hamiltonian with the Phillips--Kleinman
projection added:
\begin{equation*}
  H_\text{val} \;\longrightarrow\; H_\text{val}
  \;+\; \sum_{c \in \text{core}} (E_v - E_c)\,
  |\phi_c\rangle\langle\phi_c|\,,
\end{equation*}
or equivalently at the matrix-element level
$\Delta H_{pq}^\text{PK} = \sum_c (E_v - E_c)\,S_{pc}\,S_{cq}$ for
valence indices $p,q$ and cross-overlap $S_{pc} = \langle\phi_p\,
|\,\phi_c\rangle$, with $\{\phi_c, E_c\}$ the (external) frozen-core
orbitals and binding energies and $E_v$ a valence reference energy
(typically $E_v = 0$ for the standard absolute-PK barrier, which is
purely repulsive whenever $E_c < 0$).\\
\textbf{Variables introduced:} the indexed set
$\{\phi_c, E_c\}_{c \in \text{core}}$ of core orbital wavefunctions
and binding energies.  Structurally one bookkeeping integer (the
core orbital count) plus an external Layer-2 input for each $E_c$ and
each radial profile $\phi_c$.  These are not generated by the
framework's bare graph or projection chain; they are consumed from
Clementi--Roetti HF tabulations or from the framework's own core-shell
solution (Paper~17~\S\ref{sec:pk}'s \texttt{CoreScreening} object,
which solves the same-center core-electron problem independently and
hands its eigenvectors to the valence solver).\\
\textbf{Dimension introduced:} none.  Orthogonality is dimensionless;
the projector $\sum_c |\phi_c\rangle\langle\phi_c|$ takes the same
dimension as its argument.  The valence reference $E_v$ and the
core eigenvalues $E_c$ carry energy, but they were already part of the
spectrum-of-states ledger before the projection was applied.  In this
sense Phillips--Kleinman is in the same dimensionless class as
Drake--Swainson asymptotic subtraction (\S\ref{sec:proj_drake_swainson})
--- both are bookkeeping projections that re-organize an already-
dimensioned quantity rather than introduce a new physical scale.\\
\textbf{Transcendental signature:} ring-preserving.  No transcendental
is injected beyond what the source spectrum already carries.  In the
finite hydrogenic / Slater basis used by the composed and balanced
architectures, the cross-overlaps $S_{pc}$ are evaluated in closed
form (radial integrals via the hypergeometric Slater machinery
\texttt{geovac/hypergeometric\_slater.py}) and live in $\mathbb{Q}$ at
the angular level (Wigner~$3j$, \S\ref{sec:proj_3j}) and in the same
algebraic-extension ring as the underlying orbitals at the radial
level.  The PK weights $(E_v - E_c)$ are external Layer-2 numerical
inputs whose transcendental content is whatever the source compilation
delivers.  No new $\pi$, no new $\zeta$ value, no new Catalan-class
transcendental enters from the projection itself.\\
\textbf{Used in:} the same-center PK barrier in the composed
architecture (Paper~17~\S\ref{sec:pk}, module
\texttt{geovac/ab\_initio\_pk.py}) for LiH and BeH$^+$ at
$l_\text{max} \leq 2$; the $l$-dependent variant
(Paper~17~\S\ref{sec:pk}, channel-blind $\to$ $l = 0$ projected) which
reduces LiH $R_\text{eq}$ error from $6.9\%$ to $5.3\%$; the
cross-center PK barrier
(Paper~19~\S\ref{sec:w1c_residual}, module
\texttt{geovac/phillips\_kleinman\_cross\_center.py}) for second-row
chemistry where the heavy-atom core is frozen out via a
\texttt{FrozenCore} potential rather than carried as an explicit
electron block.\\
\textbf{Empirical anchor:} LiH balanced-coupled $R_\text{eq}$ at
$5.3\%$ error with the channel-blind $l$-dependent PK at
$l_\text{max} = 2$ (Paper~17~\S\ref{sec:pk}, Track~CD baseline; the
PK contribution is necessary --- removing it produces an unbound PES,
see Track~CB record in CLAUDE.md~\S~3); core-valence
orthogonality satisfied numerically in the explicit-core architectures
(LiH balanced-coupled at $n_\text{max}=2$ has the Li-core block
diagonalized simultaneously with the LiH-bond block; Pauli is
enforced by Slater-determinant antisymmetry in the FCI sector and the
PK projector is redundant for the explicit-core slice but
load-bearing for the frozen-core slice).\\
\textbf{Honest scope.}  Three layers of scope-honesty are required
here, all flagged for CLAUDE.md~\S~1.5 rhetoric compliance.

First, in the composed / balanced architecture as implemented, the PK
projection is \emph{essential} --- it is not redundant in second
quantization.  The Track~CB sprint (CLAUDE.md~\S~3 record, April~2026)
tested whether explicit cross-block ERIs could substitute for the PK
barrier by carrying core-valence repulsion explicitly; the result was
that adding cross-block $V_{ee}$ without simultaneously balancing
cross-center $V_{ne}$ produces a $29\%$ FCI energy error at LiH
$n_\text{max}=2$ (worst of three tested configurations, CLAUDE.md~\S~3
record).  The composed orbital basis uses different effective charges
$Z$ on different blocks, and PK is the infrastructure that handles
the corresponding orthogonality without requiring two-center one-body
integrals in a single-center hydrogenic basis.  PK is \emph{not} a
calibration choice; it is the projection that makes the composed
architecture's Hilbert-space factorization consistent with Pauli
exclusion.

Second, the same-center PK in
Paper~17~\S\ref{sec:pk}~\cite{paper17,pk1959} closes the
core-valence-orthogonality problem cleanly at the $l_\text{max} = 2$
operating point.  Six prior PK modification attempts
(channel-blind, projector, spectral, self-consistent, $R$-dependent,
$l$-dependent on the 2D solver; see the
six-row PK modifications digest in the CLAUDE.md~\S~3
failed-approaches table) all carry the same lesson:
PK provides coordinate-space exclusion that angular projectors cannot
replicate.  The $l_\text{max}$ divergence (Paper~17~\S\ref{sec:rho_collapse}
and outlook section) at higher $l_\text{max}$ is a structural feature
of the composed architecture's interaction between PK and angular
basis expansion, not a deficiency of the PK projection itself; the
2D variational solver reproduces the same drift, confirming that PK
is not the bottleneck (Paper~17 Track~BQ record).

Third, the cross-center PK barrier
(Paper~19~\S\ref{sec:w1c_residual}, May~2026) extends the same
projection to the heavy-atom-frozen-core regime
(Z $\geq 11$, [Ne] / [Ar] / [Kr] / [Xe] cores) and was tested as the
named engineering closure for the W1c-residual orthogonality wall
identified in the Phase~C chemistry sprint.  The honest verdict was
NEGATIVE on binding outcome but POSITIVE on architectural completeness:
PK cross-center reduces the W1c-residual NaH overattraction by
$14.6\%$ ($0.357$ Ha $\to$ $0.305$ Ha at standard absolute-PK
$E_v = 0$) but does \emph{not} close the wall.  The NaH PES remains
monotonically descending with no internal equilibrium minimum.  The
remaining residual ($\sim 290$~mHa, $4\times$ experimental
$D_e \approx 75$~mHa) lives in mechanisms beyond W1c + PK orthogonality:
the hydrogenic $Z_\text{orb} = 1$ basis on the Na valence side is
qualitatively wrong (it is not the actual Na 3s orbital but a generic
$Z_\text{eff} = 1$ hydrogenic function), the cross-block ERI between
H~1s (proper) and Na ``valence'' (wrong shape) inherits this error,
and a strong-PK probe at unphysical $E_v = 100$~Ha shows a
minimum-depth signature ($0.148$~Ha descent) that represents parameter
tuning rather than derivation.  The CLAUDE.md~\S~3 failed-approaches
table carries the explicit row ``Phillips-Kleinman cross-center
barrier for second-row chemistry binding (post-Track-2 sprint,
2026-05-08)'' documenting this honest negative.  The PK projection
remains a clean architectural addition; the wall is structurally
deeper than PK orthogonality alone.\\
\textbf{Structural reading.}  Phillips--Kleinman sits in the
Layer-2~$\to$~Layer-2 input class, structurally a sibling of the three
nuclear-structure projections (\S\ref{sec:proj_charge_density},
\S\ref{sec:proj_magnetization_density}, \S\ref{sec:proj_tensor_multipole}):
each takes the framework's bare single-particle output and adds a
structural correction whose load-bearing data is external to the
framework's graph and projection chain.  For the nuclear-structure
trio the external data is a QCD-internal scalar
($\langle r^2\rangle_E$, $r_Z$, $Q_N$); for Phillips--Kleinman the
external data is the set of frozen-core orbitals
$\{\phi_c, E_c\}$ derivable either from an independent atomic
Hartree--Fock solve (Clementi--Roetti) or from the framework's own
core-shell Level~3 solution.  In both cases the projection mechanism
is structurally well-defined and ring-preserving; in both cases the
input scalars / orbitals are calibration content the framework does
not autonomously derive.

Two further structural distinctions deserve naming.  First,
Phillips--Kleinman differs categorically from the three
nuclear-structure projections in its targeted Hilbert-space sector:
the nuclear-structure projections modify the
electron-nucleus or hyperfine operators by convolving with a nuclear
density profile, while Phillips--Kleinman modifies the kinetic-energy
\emph{plus} potential-energy effective single-particle Hamiltonian
on the valence electrons via an orthogonality projector.  Different
mechanism, same input class.

Second, the same-center and cross-center variants of PK are
structurally one projection in two architectures.  In the same-center
(explicit core block) architecture, the core orbitals live inside the
FCI sector and Pauli is enforced by Slater-determinant antisymmetry;
the PK barrier in Paper~17~\S\ref{sec:pk} is the residual content
that handles the kinetic-energy cost of orthogonalizing the angular
component of the valence basis against the radial core.  In the
cross-center (frozen-core) architecture, the core orbitals are
external to the FCI Hilbert space (they live in a one-body local
$Z_\text{eff}(r)$ potential), and the PK barrier (Paper~19
\S\ref{sec:w1c_residual}) is the non-local one-body operator that
re-injects core-valence orthogonality the frozen-core approximation
removed.  Same operator-form projection, two architectures, two
load-bearing roles, the same load-bearing limitation: PK gets the
projection structure right but cannot independently fix a wrong
valence basis.
